\documentclass[17pt,twoside, openany]{extbook}
\usepackage[T1]{fontenc}
\usepackage[utf8]{inputenc}
\usepackage[french]{babel}
%\usepackage{fouriernc} % on utilise la police New Century Schoolbook
\usepackage{kpfonts} % un peu plus classique
\usepackage{geometry}
\usepackage{titling}
\newcommand{\subtitle}[1]{%
  \posttitle{%
    \par\end{center}
    \begin{center}\large#1\end{center}
    \vskip0.5em}%
}
\usepackage{subfiles}
\geometry{a4paper, left=20mm, right=20mm, top=20mm, bottom=20mm}
\author{Abbé HUMBERT}
\title{Instructions sur les principales vérités de la Religion, et sur les principaux devoirs du christianisme}
\subtitle{\small{Nouvelle édition, augmentée de nouveaux chapitres et de résolutions pratiques placées à la fin de chaque chapitre, avec une courte prière; et d'une notice sur l'Auteur}}
\date{1833}


\begin{document}
    \maketitle
    \tableofcontents
    \subfile{notice.tex}
    \subfile{001-Chapitre-premier}
    \subfile{002-Chapitre-second}
    \subfile{003-Chapitre-troisieme}
    \subfile{004-Chapitre-quatrieme}
    \subfile{005-Chapitre-cinq}
    \subfile{006-Chapitre-six}
    \subfile{007-Chapitre-sept}
    \subfile{008-Chapitre}
    \subfile{009-Chapitre}
    \subfile{010-Chapitre}
    \subfile{011-Chapitre}
    \subfile{012-Chapitre}
    \subfile{013-Chapitre}
    \subfile{014-Chapitre}
    \subfile{015-Chapitre}
    \subfile{016-Chapitre}
    \subfile{017-Chapitre}
    \subfile{018-Chapitre}
    \subfile{019-Chapitre}
    \subfile{020-Chapitre}
    \subfile{021-Chapitre}
    \subfile{022-Chapitre}
\end{document}
